\documentclass[11pt,a4paper]{article}
\usepackage[margin=1in]{geometry}
\usepackage{amsmath,amssymb,amsthm}
\usepackage{algorithm}
\usepackage{algpseudocode}
\usepackage{booktabs}
\usepackage{hyperref}

\theoremstyle{definition}
\newtheorem{definition}{Definition}[section]
\newtheorem{theorem}{Theorem}[section]
\newtheorem{lemma}[theorem]{Lemma}
\newtheorem{remark}{Remark}[section]

\title{\textbf{Rigorous Mathematical Specification, Layer-Wise Trust Ratio Dynamics, and Convergence of LARS and LAMB Optimizers}}
\author{Team Scaibu \\ \small Email: \texttt{jaydeep.9664920749@gmail.com} \\ \small Architecture and Core Engine Team}
\date{October 2026}

\begin{document}
\maketitle

\begin{abstract}
This document presents the complete theoretical formulation and algorithmic specification of Layer-wise Adaptive Rate Scaling (LARS) and Layer-wise Adaptive Moments for Batch training (LAMB). Designed to stabilize neural network training under massive distributed batch sizes (e.g., $32\text{k}$ to $64\text{k}$), both algorithms compute layer-specific trust ratios $r_t = \phi \frac{\|\theta\|_2}{\|u_t\|_2}$ to prevent exploding or vanishing updates across layers with disparate gradient-to-weight norm ratios. We establish formal layer-wise norm bounds, provide the complete algorithmic pseudo-code, derive computational complexities, demonstrate a worked numerical example, and discuss edge-case numerical considerations.
\end{abstract}

\section{Notation and Mathematical Preliminaries}

Let $\theta \in \mathbb{R}^P$ denote the parameter vector belonging to a specific network layer or parameter tensor. Optimization proceeds across iterations $t \in \{1, \dots, T\}$.

\begin{itemize}
    \item $g_t \triangleq \nabla_\theta \mathcal{L}(\theta_t) \in \mathbb{R}^P$: Stochastic gradient vector for the layer.
    \item $\|\theta\|_2 \triangleq \sqrt{\sum_{i=1}^P \theta_i^2}$: Euclidean norm of the layer parameter tensor.
    \item $m_t, v_t \in \mathbb{R}^P$: First and second moment exponential moving averages (for LAMB).
    \item $u_t \in \mathbb{R}^P$: Unscaled raw layer update vector.
    \item $r_t \in \mathbb{R}_{> 0}$: Layer-wise trust ratio scaling factor.
    \item $\phi > 0$: Trust coefficient hyperparameter (typically $\phi = 1.0$).
    \item $\lambda \ge 0$: Weight decay coefficient.
    \item $\eta > 0$: Global learning rate.
\end{itemize}

\begin{table}[h]
\centering
\caption{Summary of Mathematical Symbols for LARS / LAMB}
\begin{tabular}{lll}
\toprule
\textbf{Symbol} & \textbf{Type / Domain} & \textbf{Description} \\
\midrule
$\theta_t$ & $\mathbb{R}^P$ & Layer parameter vector \\
$g_t$ & $\mathbb{R}^P$ & Layer gradient vector \\
$u_t$ & $\mathbb{R}^P$ & Layer-wise raw update vector \\
$r_t$ & $\mathbb{R}_{> 0}$ & Dynamic trust ratio $\phi \frac{\|\theta\|_2}{\|u_t\|_2}$ \\
$\phi$ & $\mathbb{R}_{> 0}$ & Trust coefficient ($1.0$) \\
$\lambda$ & $\mathbb{R}_{\ge 0}$ & Weight decay parameter \\
$\eta$ & $\mathbb{R}_{> 0}$ & Global base learning rate \\
$\beta_1, \beta_2$ & $[0, 1)$ & Moment decay parameters ($0.9, 0.999$ for LAMB) \\
\bottomrule
\end{tabular}
\end{table}

\section{Problem Definition and Core Concepts}

\subsection{Intuition and Motivation}
When scaling batch sizes from small minibatches (e.g., 256) to massive parallel batches (e.g., 32,768), linear learning rate scaling $\eta \propto B$ causes training divergence during early iterations. Deep investigation revealed that the ratio of weight norm to gradient norm $\|\theta^{(l)}\|_2 / \|g^{(l)}\|_2$ varies by orders of magnitude across different network layers (e.g., up to $1000\times$ between input embeddings and deep feed-forward projections). Unscaled global updates cause unstable catastrophic updates in layers with small weight norms while starving layers with large weight norms. LARS and LAMB normalize update magnitudes on a per-layer basis using the trust ratio $r_t$.

\subsection{Formal Mathematical Definition}
\begin{definition}[LARS vs. LAMB Update Formulations]
For a parameter tensor $\theta \in \mathbb{R}^P$:
\begin{itemize}
    \item \textbf{LARS (Stochastic Gradient)}:
    \begin{equation}
    u_t = g_t + \lambda \theta_t
    \end{equation}
    \item \textbf{LAMB (Adaptive Moments)}:
    \begin{align}
    m_t &= \beta_1 m_{t-1} + (1 - \beta_1) g_t, \quad v_t = \beta_2 v_{t-1} + (1 - \beta_2) g_t^2 \\
    \hat{m}_t &= \frac{m_t}{1 - \beta_1^t}, \quad \hat{v}_t = \frac{v_t}{1 - \beta_2^t} \\
    u_t &= \frac{\hat{m}_t}{\sqrt{\hat{v}_t} + \epsilon} + \lambda \theta_t
    \end{align}
\end{itemize}
The layer trust ratio $r_t$ and parameter update are computed as:
\begin{align}
r_t &= \begin{cases} \phi \frac{\|\theta_t\|_2}{\|u_t\|_2} & \text{if } \|\theta_t\|_2 > 0 \text{ and } \|u_t\|_2 > 0 \\ 1.0 & \text{otherwise} \end{cases} \label{eq:trust_ratio} \\
\theta_{t+1} &= \theta_t - \eta \cdot r_t \cdot u_t \label{eq:lars_lamb_update}
\end{align}
\end{definition}

\section{Algorithmic Specification}

The unified execution flow from \texttt{NnAlgoLarsLambOptimizer} is presented below.

\begin{algorithm}[H]
\caption{Layer-Wise Adaptive Optimization (LARS / LAMB)}
\begin{algorithmic}[1]
\Require Parameter tensor $\theta \in \mathbb{R}^P$, gradient $g \in \mathbb{R}^P$, mode $\in \{\text{lars}, \text{lamb}\}$
\Require Step $t \ge 1$, base rate $\eta > 0$, weight decay $\lambda \ge 0$, trust coefficient $\phi > 0$, stability $\epsilon > 0$
\Require Optional moment buffers $m, v \in \mathbb{R}^P$ (for LAMB)
\Ensure Updated parameter vector $\theta_{\text{next}}$, scalar trust ratio $r_t$, updated moments $(m_{\text{next}}, v_{\text{next}})$
\State $P \gets \text{length}(\theta)$
\If{$P = 0$ \textbf{or} $\text{length}(g) \ne P$}
    \State \textbf{throw} PreconditionFailureException(``Buffer length mismatch or zero length'')
\EndIf
\State $w_{\text{norm}} \gets \sqrt{\sum_{i=1}^P \theta[i]^2}$
\State $u \gets \text{new Array of size } P$
\State $m_{\text{next}} \gets \text{new Array of size } P$, \quad $v_{\text{next}} \gets \text{new Array of size } P$
\If{$\text{mode} = \text{lamb}$}
    \State $b_1 \gets 1.0 - \beta_1^t, \quad b_2 \gets 1.0 - \beta_2^t$
    \For{$i \gets 1 \textbf{ to } P$}
        \State $m_{\text{next}}[i] \gets \beta_1 \cdot m[i] + (1.0 - \beta_1) \cdot g[i]$
        \State $v_{\text{next}}[i] \gets \beta_2 \cdot v[i] + (1.0 - \beta_2) \cdot (g[i])^2$
        \State $\hat{m} \gets m_{\text{next}}[i] / b_1, \quad \hat{v} \gets v_{\text{next}}[i] / b_2$
        \State $u[i] \gets \left(\frac{\hat{m}}{\sqrt{\hat{v}} + \epsilon}\right) + \lambda \cdot \theta[i]$
    \EndFor
\Else
    \For{$i \gets 1 \textbf{ to } P$}
        \State $u[i] \gets g[i] + \lambda \cdot \theta[i]$
    \EndFor
\EndIf
\State $u_{\text{norm}} \gets \sqrt{\sum_{i=1}^P u[i]^2}$
\If{$w_{\text{norm}} > 0$ \textbf{and} $u_{\text{norm}} > 0$}
    \State $r_t \gets \phi \cdot (w_{\text{norm}} / u_{\text{norm}})$
\Else
    \State $r_t \gets 1.0$
\EndIf
\State $\theta_{\text{next}} \gets \text{new Array of size } P$
\For{$i \gets 1 \textbf{ to } P$}
    \State $\theta_{\text{next}}[i] \gets \theta[i] - \eta \cdot r_t \cdot u[i]$
\EndFor
\State \Return $\{\text{updated\_parameters}: \theta_{\text{next}}, \text{trust\_ratio}: r_t, \text{updated\_exp\_avg}: m_{\text{next}}, \text{updated\_exp\_avg\_sq}: v_{\text{next}}\}$
\end{algorithmic}
\end{algorithm}

\section{Theoretical Guarantees and Step-Size Bounds}

\begin{theorem}[Strict Layer-Wise Relative Step-Size Invariance]
For any layer tensor $\theta_t$ with $\|\theta_t\|_2 > 0$ and non-zero raw update $\|u_t\|_2 > 0$, the effective Euclidean norm of the parameter update $\Delta \theta_t = \theta_{t+1} - \theta_t$ satisfies:
\begin{equation}
\|\Delta \theta_t\|_2 = \eta \cdot \phi \cdot \|\theta_t\|_2
\end{equation}
Consequently, the relative parameter update ratio $\frac{\|\Delta \theta_t\|_2}{\|\theta_t\|_2} = \eta \phi$ is invariant to gradient magnitude scaling and layer width.
\end{theorem}

\begin{proof}
By Equation~\eqref{eq:lars_lamb_update}, $\Delta \theta_t = -\eta \cdot r_t \cdot u_t$. Taking the Euclidean norm:
\begin{equation}
\|\Delta \theta_t\|_2 = \|-\eta r_t u_t\|_2 = \eta r_t \|u_t\|_2
\end{equation}
Substituting $r_t = \phi \frac{\|\theta_t\|_2}{\|u_t\|_2}$:
\begin{equation}
\|\Delta \theta_t\|_2 = \eta \left(\phi \frac{\|\theta_t\|_2}{\|u_t\|_2}\right) \|u_t\|_2 = \eta \phi \|\theta_t\|_2
\end{equation}
Dividing both sides by $\|\theta_t\|_2$ gives exactly $\frac{\|\Delta \theta_t\|_2}{\|\theta_t\|_2} = \eta \phi$.
\end{proof}

\subsection{Limitations}
\begin{itemize}
    \item \textbf{Exclusion for 1D Parameters}: Trust ratio scaling should be bypassed ($r_t = 1.0$) for 1D biases and layer-norm scales, where norm statistics across 1D slices produce noisy ratios.
\end{itemize}

\section{Complexity Analysis}

\subsection{Time Complexity}
\begin{itemize}
    \item \textbf{Reduction Passes}: Two $L_2$ norm reductions across $P$ elements ($\|\theta\|_2$ and $\|u\|_2$), requiring $2P$ multiply-accumulate operations ($\mathcal{O}(P)$).
    \item \textbf{Element Updates}: In LAMB, moment calculations require $\mathcal{O}(P)$ operations.
    \item \textbf{Total Time}: $\Theta(P)$ FLOPs per layer per step.
\end{itemize}

\subsection{Space Complexity}
\begin{itemize}
    \item \textbf{Auxiliary Memory}: In LARS, $\mathcal{O}(1)$ additional memory beyond updates. In LAMB, $2P$ floats for moment buffers $m, v$ ($\Theta(P)$ auxiliary space).
\end{itemize}

\section{Worked Numerical Example}

Consider a $2$-element parameter vector $\theta = [3.0, 4.0]^T$, gradient $g = [1.0, 2.0]^T$, mode $=$ \texttt{"lars"}, $\eta = 0.01$, $\lambda = 0.1$, $\phi = 1.0$.

\begin{enumerate}
    \item \textbf{Parameter Norm}:
    \begin{equation}
    w_{\text{norm}} = \sqrt{3.0^2 + 4.0^2} = \sqrt{9.0 + 16.0} = \sqrt{25.0} = 5.0
    \end{equation}
    \item \textbf{Raw Updates}:
    \begin{align}
    u_1 &= g_1 + \lambda \theta_1 = 1.0 + 0.1(3.0) = 1.30 \\
    u_2 &= g_2 + \lambda \theta_2 = 2.0 + 0.1(4.0) = 2.40
    \end{align}
    \item \textbf{Update Norm}:
    \begin{equation}
    u_{\text{norm}} = \sqrt{1.3^2 + 2.4^2} = \sqrt{1.69 + 5.76} = \sqrt{7.45} \approx 2.7294688
    \end{equation}
    \item \textbf{Trust Ratio}:
    \begin{equation}
    r_t = \phi \frac{w_{\text{norm}}}{u_{\text{norm}}} = 1.0 \times \frac{5.0}{2.7294688} \approx 1.8318583
    \end{equation}
    \item \textbf{Parameter Updates}:
    \begin{align}
    \theta_{\text{next}, 1} &= 3.0 - (0.01)(1.8318583)(1.30) = 3.0 - 0.02381416 = 2.97618584 \\
    \theta_{\text{next}, 2} &= 4.0 - (0.01)(1.8318583)(2.40) = 4.0 - 0.04396460 = 3.95603540
    \end{align}
\end{enumerate}

\section{Numerical Considerations and Edge Cases}

\begin{itemize}
    \item \textbf{Zero Weight Initialization}: If a parameter tensor is initialized to all zeros ($\|\theta\|_2 = 0$), $r_t$ would evaluate to $0.0$, freezing the layer. The algorithm explicitly guards $w_{\text{norm}} > 0$ and falls back to $r_t = 1.0$.
    \item \textbf{Distributed Reductions}: In tensor-parallel and data-parallel regimes, $\|\theta\|_2$ and $\|u\|_2$ must be reduced globally across all sharded tensor partitions using \texttt{AllReduce} before computing $r_t$.
\end{itemize}

\begin{thebibliography}{9}
\bibitem{you2017} Y.~You, I.~Gitman, and B.~Ginsburg, ``Large batch training of convolutional networks,'' \emph{arXiv preprint arXiv:1708.03888}, 2017.
\bibitem{you2019} Y.~You et al., ``Large batch optimization for deep learning: Training BERT in 76 minutes,'' in \emph{International Conference on Learning Representations (ICLR)}, 2020.
\bibitem{ginsburg2019} B.~Ginsburg et al., ``Stochastic Gradient Methods with Layer-wise Adaptive Moments for Training of Deep Networks,'' \emph{arXiv preprint arXiv:1904.00962}, 2019.
\end{thebibliography}

\end{document}
